\chapter{异构CAD特征差异分析与转换需求}
\label{chap:2}
\chapterguide{按参数、空间、组织和依赖等维度归纳问题，并逐类落实到五类特征。此章分析“差异是什么”，不提前展开完整算法。}

\section{研究对象与差异分析框架}
\label{sec:2-1}

\subsection{三系统参数调查的对象与方法}
\label{subsec:2-1-1}
\outlinehint{以SolidWorks、Creo、NX的特征定义、读取参数和创建条件为调查对象，记录系统版本、数据来源和可访问范围。}

\subsection{九类异构问题的分类与对应关系}
\label{subsec:2-1-2}
\outlinehint{保留九类分析框架：量纲、方向、参数含义、对象组织、约束依赖、引用历史、几何拓扑、功能范围、执行环境；建立分类与五类特征的交叉表。}

\section{草图与基准面的表达差异}
\label{sec:2-2}

\subsection{草图坐标、图元与圆弧表达差异}
\label{subsec:2-2-1}
\outlinehint{比较局部坐标与世界坐标、圆弧范围、构造线以及椭圆读取和创建范围；通过最小案例说明直接复制数值的风险。}

\subsection{草图约束、尺寸与外部关联差异}
\label{subsec:2-2-2}
\outlinehint{比较约束参与对象、起终点和圆心等作用位置，区分驱动尺寸、几何关系和对外部几何的关联。}

\subsection{基准面构造方法与参考条件差异}
\label{subsec:2-2-3}
\outlinehint{比较偏置、角度、中间面、过点过线等构造和参考组合，说明相同平面位置不等于相同构造依赖。}

\section{拉伸、旋转与扫掠的参数差异}
\label{sec:2-3}

\subsection{拉伸方向、深度与终止条件差异}
\label{subsec:2-3-1}
\outlinehint{对比单向、双向、对称、到对象与到延伸面等含义，使用总深度60 mm与单侧30 mm说明换算需求。}

\subsection{旋转轴、转动方向与角度范围差异}
\label{subsec:2-3-2}
\outlinehint{区分轴线本身的方向和特征的正转方向，比较两侧角度的正负约定以及参考终止可能产生的互补角。}

\subsection{扫掠截面、路径与方向控制差异}
\label{subsec:2-3-3}
\outlinehint{比较独立草图、内嵌截面、参数圆截面，分析路径起点、曲线段顺序、截面放置和沿路径方向控制的共同影响。}

\section{跨特征的共性差异}
\label{sec:2-4}

\subsection{数值量纲、模板单位与薄壁材料位置差异}
\label{subsec:2-4-1}
\outlinehint{区分长度、角度与无量纲参数；比较厚度加方向与两侧偏置，说明总厚度相同仍可能形成不同的材料位置。}

\subsection{引用身份、建模历史与拓扑划分差异}
\label{subsec:2-4-2}
\outlinehint{说明原生编号不可跨系统复用，分析父特征归属变化、边分段与接缝对引用恢复和几何比较的影响。}

\subsection{目标功能范围与构造组织差异}
\label{subsec:2-4-3}
\outlinehint{归纳独立草图与内嵌截面、辅助基准面、子类型不对称、参数未使用和原生执行条件的差异。}

\section{转换需求与评价目标归纳}
\label{sec:2-5}

\subsection{公共表达与目标重建需求}
\label{subsec:2-5-1}
\outlinehint{从差异推导统一量纲、明确方向、带角色的引用、约束作用位置及可解释组合重建等需求，形成可追踪的需求表。}

\subsection{能力边界与结果评价需求}
\label{subsec:2-5-2}
\outlinehint{区分结构可表达、源端能读取、目标能创建和实验证明四个层次；提出结构、参数、几何和编辑再生的分别检查要求。}

\section{本章小结}
\label{sec:2-6}

\outlinehint{归纳本章能够支持的认识，并说明与下一章的衔接；不重复列举所有接口或扩展未证明结论。}
